\documentclass[12pt]{article}
\usepackage[T1]{fontenc}
\usepackage{lmodern}
\usepackage{amsmath}
\usepackage{amssymb}
\usepackage{amsthm}
\usepackage[a4paper,margin=1in]{geometry}
\usepackage{microtype}
\usepackage{needspace}
\usepackage{titling}
\usepackage{fancyhdr}
\usepackage[colorlinks=true,linkcolor=blue,citecolor=blue,urlcolor=blue]{hyperref}
\setlength{\droptitle}{-4.75em}
\pretitle{\begin{center}\Large\bfseries}
\posttitle{\par\end{center}\vspace{-0.5em}}
\preauthor{\begin{center}\large}
\postauthor{\par\vspace{0.5em}
    {\fontsize{11}{13}\selectfont\itshape Communicated and edited by Boming Jia\par}
    \end{center}\vspace{-1.75em}}
\predate{}
\postdate{}
\setlength{\headheight}{15pt}
\pagestyle{fancy}
\fancyhf{}
\fancyhead[L]{\itshape On the minimal orbital scheme}
\fancyfoot[C]{\thepage}
\renewcommand{\headrulewidth}{0pt}
\setlength{\parindent}{1.5em}
\setlength{\parskip}{0pt}
\newtheorem{theorem}{Theorem}[section]
\newtheorem{lemma}[theorem]{Lemma}
\newtheorem{proposition}[theorem]{Proposition}
\theoremstyle{definition}
\newtheorem{example}[theorem]{Example}
\DeclareMathOperator{\Sym}{Sym}
\DeclareMathOperator{\height}{ht}
\DeclareMathOperator{\inid}{in}
\DeclareMathOperator{\LT}{LT}
\DeclareMathOperator{\Hilb}{Hilb}
\DeclareMathOperator{\Ad}{Ad}
\newcommand{\C}{\mathbb C}
\newcommand{\g}{\mathfrak g}
\newcommand{\h}{\mathfrak h}
\newcommand{\n}{\mathfrak n^{+}}
\newcommand{\nm}{\mathfrak n^-}
\newcommand{\omin}{\mathcal O_{\min}}
\newcommand{\oc}{\overline{\mathcal O_{\min}}}
\newcommand{\PP}{\mathbb P}
\title{On the Minimal Orbital Scheme $\oc\cap\n$}
\author{GPT-6 Astra}
\date{}
\hypersetup{pdftitle={On the Minimal Orbital Scheme Omin-bar intersect n+},pdfauthor={GPT-6 Astra},pdfsubject={Scheme-theoretic intersections, quadratic Groebner bases, and Springer fibers}}
\begin{document}
\maketitle
\thispagestyle{plain}
\begin{abstract}
    Let $\mathfrak g$ be a complex simple Lie algebra, let $\mathfrak n^{+}$ be the nilradical of a Borel subalgebra, and let $\mathcal O_{\min}$ be the minimal nonzero nilpotent orbit. We prove that $\overline{\mathcal O_{\min}}\cap\mathfrak n^{+}$ is generically reduced in every type. In simply laced type, it is reduced and Cohen--Macaulay. In non-simply laced type, it is nonreduced and, being generically reduced, fails Serre's condition $(S_1)$. In simply laced type, a quadratic Gr\"obner basis degenerates its coordinate ring to the Stanley--Reisner ring of the positive-root order complex. As a consequence, minimal scheme-theoretic Springer fibers are generically reduced in every type, reduced and Cohen--Macaulay in simply laced type, and nonreduced in non-simply laced type.
\end{abstract}

\section{Introduction}

Let $G$ be a connected complex simple algebraic group with Lie algebra $\g$, and let $B\subset G$ be a Borel subgroup. Let $\n$ be the nilradical of $\operatorname{Lie}B$, let $\omin$ be the minimal nonzero nilpotent orbit, and let $\Phi$ be the root system of $G$. All intersections below are scheme-theoretic.

Jia and Li proved that $\oc\cap\n$ is reduced and Cohen--Macaulay in type $A$ by a squarefree Gr\"obner degeneration \cite[Theorem~1.4]{JL25}.

\begin{theorem}\label{thm:main}
    The scheme $\overline{\mathcal O_{\min}}\cap\mathfrak n^{+}$ is generically reduced.
    \begin{enumerate}
        \item[(i)] If $\Phi$ is simply laced, then $\overline{\mathcal O_{\min}}\cap\mathfrak n^{+}$ is reduced and Cohen--Macaulay.
        \item[(ii)] If $\Phi$ is not simply laced, then $\overline{\mathcal O_{\min}}\cap\mathfrak n^{+}$ is nonreduced and, being generically reduced, fails Serre's condition $(S_1)$.
    \end{enumerate}
\end{theorem}

Choose a maximal torus $T\subset B$, with positive roots $\Phi^+$. Write $\alpha\leq\beta$ when $\beta-\alpha$ is a nonnegative integral combination of simple roots. Let $I\subset\C[\n]$ define $\oc\cap\n$. Enumerate the positive roots by increasing height, with the simple roots first, and use the variable order
$$
    x_{\beta_1}\prec\cdots\prec x_{\beta_N},\qquad N=|\Phi^+|.
$$
For graded reverse lexicographic order $\prec$, we prove in simply laced type that $I$ has a quadratic Gr\"obner basis with
$$
    \inid_\prec I=(x_\alpha x_\beta\mid\alpha,\beta\in\Phi^+\text{ are incomparable}).
$$
\par\medskip
\noindent\textbf{Human Editor's Remark.}
This paper was generated by GPT-6 Astra and revised together with the Human Editor (HE), Boming Jia. HE has found no mathematical errors in this paper. If any errors are found later, HE will take full responsibility. If the proof is indeed correct, HE does not want any credit for HE's negligible mathematical contribution.

\section{Root coordinates and generic reducedness}

Let $\h=\operatorname{Lie}T$ and let $\nm$ be the opposite nilradical, so that
$$
    \g=\nm\oplus\h\oplus\n.
$$
Write $r=\operatorname{rank}\g$ and let $\theta$ be the highest root. Let $\alpha_1,\ldots,\alpha_r$ be the simple roots and set
$$
    \height\left(\sum_i c_i\alpha_i\right)=\sum_i c_i.
$$
We use Bourbaki's conventions \cite[Chapter~VI]{Bou02}.

For each root $\alpha$, choose a nonzero root vector $e_\alpha$. The coordinate $x_\alpha$ is defined by
$$
    x_\alpha(e_\beta)=\delta_{\alpha\beta},\qquad x_\alpha|_\h=0.
$$
Its $T$-weight is $-\alpha$, and
$$
    \C[\n]=\C[x_\alpha\mid\alpha\in\Phi^+].
$$

Enumerate $\Phi^+$ as $\beta_1,\ldots,\beta_N$ by increasing height. At equal height, use decreasing lexicographic order on the coefficient tuples in the ordered basis $\alpha_1,\ldots,\alpha_r$. Set
$$
    x_i=x_{\beta_i},\qquad \beta_i=\alpha_i\quad(1\leq i\leq r).
$$
Thus the simple-root coordinates come first. Throughout, $\prec$ denotes graded reverse lexicographic order with
$$
    x_1\prec\cdots\prec x_N.
$$
Explicitly, among distinct monomials of the same total degree, the larger monomial is the one with the smaller exponent at the least index where their exponent vectors differ.

\Needspace{9\baselineskip}
\begin{lemma}\label{lem:generic}
    The scheme $\oc\cap\n$ is generically reduced. At every nonzero closed point $x\in\oc\cap\n$, one has
    $$
        \dim T_x(\oc\cap\n)=\dim\omin-\dim(B\cdot x).
    $$
\end{lemma}

\begin{proof}
    Put $d=\frac12\dim\omin$. By Joseph's description of minimal orbital varieties \cite[\S\S1.2, 2.4]{Jos98}, the closures $\overline{B\cdot e_{\alpha_i}}$ for long simple roots $\alpha_i$ are the irreducible components of $(\oc\cap\n)_{\mathrm{red}}$, each of dimension $d$.

    Let $\kappa$ be the Killing form and $\omega$ the Kirillov--Kostant symplectic form on $\omin$. Write $\mathfrak b=\operatorname{Lie}B$. At a nonzero closed point $x\in\oc\cap\n$, one has
    $$
        \omega_x([a,x],v)=\kappa(a,v),\qquad a\in\mathfrak b,\quad v\in T_x\omin.
    $$
    Since $\omin$ is open in $\oc$ and $\mathfrak b^{\perp_\kappa}=\n$, tangent spaces to the scheme-theoretic intersection satisfy
    \begin{align*}
        T_x(\oc\cap\n)
        &=T_x\oc\cap\n
        =\bigl(T_x(B\cdot x)\bigr)^{\perp_\omega},\\
        \dim T_x(\oc\cap\n)
        &=2d-\dim(B\cdot x).
    \end{align*}
    Along each dense orbit $B\cdot e_{\alpha_i}$, the local dimension and the tangent-space dimension are both $d$. Thus $\oc\cap\n$ is smooth along these orbits and is generically reduced.
\end{proof}

The minimal orbit closure is the affine cone over $G\cdot[e_\theta]\subset\PP(\g)$. If $P_\theta$ is the stabilizer of $[e_\theta]$, then
$$
    \PP(\oc)=G/P_\theta\subset\PP(\g),
$$
where the embedding is defined by the very ample line bundle associated with $\theta$.

Let $V_{2\theta}$ be the irreducible $\g$-module of highest weight $2\theta$. By Kostant's quadratic generation theorem, as proved in \cite[Chapter~III, \S2, Theorem~1]{Gar82}, the defining ideal of $\oc$ is generated by the kernel of
$$
    \Sym^2(\g^*)\longrightarrow V_{2\theta}^*.
$$
Consequently,
\begin{equation}\label{eq:restricted}
    I=\left(q|_{\n}\ \middle|\ q\in
    \ker\bigl(\Sym^2(\g^*)\longrightarrow V_{2\theta}^*\bigr)\right).
\end{equation}
The linear equations of $\n\subset\g$ are the negative-root coordinates and the elements of $\h^*$, extended by zero on the root spaces. The resulting coordinate ring is
\begin{equation}\label{eq:quotient}
    \C[\oc\cap\n]
    =\C[\oc]/(\h^*,x_{-\alpha}\mid\alpha\in\Phi^+).
\end{equation}
\Needspace{8\baselineskip}
\begin{lemma}\label{lem:short}
    If $\alpha$ is a short simple root, then
    $$
        x_\alpha^2\in I,\qquad x_\alpha\notin I.
    $$
    If $\alpha$ is adjacent to a long simple root $\beta$, put $\gamma=s_\alpha(\beta)$. Then the class of $x_\alpha$ is a nonzero nilpotent in $\C[\oc\cap\n][x_\gamma^{-1}]$. In particular, in every non-simply laced type, $(\oc\cap\n)\setminus\{0\}$ is nonreduced.
\end{lemma}

\begin{proof}
    Since $I$ is generated by homogeneous quadrics, it contains no nonzero linear form. Thus $x_\alpha\notin I$.

    The $T$-fixed points of $\PP(\oc)$ correspond to the Weyl conjugates of $\theta$, all of which are long. Thus $e_\alpha\notin\oc$.

    By quadratic generation, there is a quadratic defining equation $q$ of $\oc$ with $q(e_\alpha)\neq0$. The space of these equations is $T$-stable. Only the component of weight $-2\alpha$ can be nonzero at $e_\alpha$, so we may choose $q$ of that weight. After restriction to $\n$, every monomial in $q$ is of the form $x_\rho x_\sigma$ with
    $$
        \rho+\sigma=2\alpha,\qquad\rho,\sigma\in\Phi^+.
    $$
    Since $\alpha$ is simple, the only possibility is $\rho=\sigma=\alpha$. Hence
    $$
        q|_{\n}=q(e_\alpha)x_\alpha^2\in I.
    $$

    Suppose that $\alpha$ is adjacent to a long simple root $\beta$. Set
    $$
        m=-\langle\beta,\alpha^\vee\rangle\in\{2,3\},
        \qquad \gamma=s_\alpha(\beta)=\beta+m\alpha.
    $$
    The root string through $\beta$ contains $\beta,\beta+\alpha,\ldots,\beta+m\alpha$. Thus $\gamma$ is a positive long root and $\gamma-\alpha$ is a root. Consequently $e_\gamma\in\omin$ and
    $$
        [e_{\alpha-\gamma},e_\gamma]=c e_\alpha,\qquad c\ne0.
    $$
    Over the dual numbers $\C[\varepsilon]/(\varepsilon^2)$, we have
    $$
        \Ad\bigl(\exp(\varepsilon e_{\alpha-\gamma})\bigr)e_\gamma
        =e_\gamma+c\varepsilon e_\alpha.
    $$
    This point lies in $\oc\cap\n$, and $x_\gamma$ takes the value $1$. Thus evaluation defines a homomorphism
    \[
        \C[\oc\cap\n][x_\gamma^{-1}]\longrightarrow\C[\varepsilon]/(\varepsilon^2),
        \qquad x_\alpha\longmapsto c\varepsilon\ne0.
    \]
    Hence the class of $x_\alpha$ remains nonzero on the principal open set $x_\gamma\ne0$, which excludes the origin. Every irreducible non-simply laced Dynkin diagram has an adjacent short--long pair, proving the last assertion.
\end{proof}

\section{The quadratic Gr\"obner basis}

For the Gr\"obner basis construction, assume that $\Phi$ is simply laced.

Let $W_\theta$ be the stabilizer of $\theta$ in the Weyl group $W$. For $w\in W/W_\theta$, let
$$
    X_w=\overline{B\cdot w[e_\theta]}\subset\PP(\oc).
$$
We identify the coset $w$ with the root $w\theta$ and write
$$
    X_\alpha=\overline{B\cdot[e_\alpha]}.
$$
Since $\n$ is $B$-stable, $X_\alpha\subset\PP(\n)$ for every $\alpha\in\Phi^+$.

\Needspace{8\baselineskip}
\begin{lemma}\label{lem:rootorder}
    For $\alpha,\beta\in\Phi^+$, one has
    $$
        X_\beta\subset X_\alpha\quad\Longleftrightarrow\quad\alpha\leq\beta.
    $$
    These conditions are also equivalent to $x_\beta$ not vanishing identically on the affine cone over $X_\alpha$.
\end{lemma}

\begin{proof}
    Normalize the invariant inner product so that every root has squared length~$2$. Suppose $\alpha<\beta$ and write
    $$
        \beta-\alpha=\sum_i n_i\alpha_i,\qquad n_i\geq0.
    $$
    Since
    $$
        (\beta-\alpha,\beta)=2-(\alpha,\beta)>0,
    $$
    there is an $i$ with $n_i>0$ and $(\alpha_i,\beta)>0$. The root $\beta$ is not simple, so $(\alpha_i,\beta)=1$. By the root-string property, $\beta-\alpha_i\in\Phi^+$, and $\alpha\leq\beta-\alpha_i$. Induction constructs a chain from $\alpha$ to $\beta$ in which each successive difference is a simple root.

    If $\gamma$ and $\gamma+\alpha_i$ are roots, then $\gamma+2\alpha_i$ is not a root. Hence
    $$
        \Ad(\exp(t e_{\alpha_i}))e_\gamma
        =e_\gamma+c t e_{\gamma+\alpha_i},\qquad c\neq0.
    $$
    Taking the projective limit, we obtain $[e_{\gamma+\alpha_i}]\in X_\gamma$. Since $X_\gamma$ is closed and $B$-stable, it contains $X_{\gamma+\alpha_i}$. Applying this along the chain proves $X_\beta\subset X_\alpha$.

    Conversely, the linear span of the affine cone over $X_\alpha$ lies in the sum of the root spaces with weights in $\alpha+\sum_i\mathbb Z_{\geq0}\alpha_i$. Thus $x_\beta$ vanishes on this cone unless $\alpha\leq\beta$. If $\alpha\leq\beta$, then $e_\beta$ lies in the cone and $x_\beta(e_\beta)=1$. This proves both the converse and the coordinate criterion.
\end{proof}

\Needspace{7\baselineskip}
\begin{lemma}\label{lem:smt}
    In simply laced type, for $\alpha\in\Phi^+$ and $d\geq1$, the monomials $x_{\gamma_1}\cdots x_{\gamma_d}$ with $\alpha\leq\gamma_1\leq\cdots\leq\gamma_d$ in $\Phi^+$ form a basis of the degree-$d$ coordinate space of the affine cone over $X_\alpha$.
\end{lemma}

\begin{proof}
    Let $\mathcal L$ be the restriction of $\mathcal O_{\PP(\g)}(1)$ to $G/P_\theta$, so that $H^0(G/P_\theta,\mathcal L)=\g^*$. Passing to the simply connected cover preserves this embedding and $\mathcal L$, allowing us to apply standard monomial theory in its simply connected formulation.

    We use the Lakshmibai--Seshadri paths of shape $\theta$ \cite[\S2]{Lit94}. For a path $\pi=(\tau_1,\ldots,\tau_s;a_0,\ldots,a_s)$, write $i(\pi)=\tau_1$ and $t(\pi)=\tau_s$ for its first and last directions in $W/W_\theta$. Its path vector $p_\pi\in\g^*$ has weight $-\pi(1)$ \cite[Theorem~3]{Lit98}. For paths of this common shape, the defining-chain criterion \cite[\S11, Corollary~3 and Theorem~8]{Lit96} says that $p_{\pi_1}\cdots p_{\pi_d}$ is standard precisely when
    $$
        t(\pi_j)\geq i(\pi_{j+1}),\qquad 1\leq j<d,
    $$
    and is standard on $X_w$ when also $i(\pi_1)\leq w$. These inequalities are in Bruhat order.

    The standard products form a basis of $H^0(G/P_\theta,\mathcal L^{\otimes d})$, and those standard on $X_w$ restrict to a basis of $H^0(X_w,\mathcal L^{\otimes d}|_{X_w})$ \cite[Theorems~4 and~7(ii)]{Lit98}. Every other standard product vanishes there because its first factor does \cite[Theorem~6]{Lit98}. Since the path vectors are degree-one coordinates, $\Sym^d(\g^*)$ maps onto both section spaces. The kernels are the homogeneous equations of the respective affine cones, so the standard products also give their homogeneous coordinate bases.

    In simply laced type, $W\cdot\theta=\Phi$, so every root $\gamma$ gives a straight L--S path $\pi_\gamma(s)=s\gamma$. Since the root weight spaces are one-dimensional, the path-vector basis theorem \cite[Theorem~3]{Lit98} shows that the remaining $r$ paths end at zero and their vectors form a basis of $\h^*$. We normalize $p_{\pi_\gamma}=x_\gamma$. Let $\mathcal B_d$ be the resulting standard monomial basis of $\C[\oc]_d$.

    For a straight path, the first and last directions coincide. Lemma~\ref{lem:rootorder} identifies Bruhat order on positive-root cosets with the reverse of root order. Hence the standard products involving only positive-root coordinates are exactly the chain monomials
    \begin{equation}\label{eq:chains}
        x_{\gamma_1}\cdots x_{\gamma_d},\qquad
        \gamma_1\leq\cdots\leq\gamma_d,\quad \gamma_j\in\Phi^+.
    \end{equation}
    Such a product is standard on $X_\alpha$ precisely when $\alpha\leq\gamma_1$. Finally, $X_\alpha\subset\PP(\n)$, so every monomial containing a negative-root coordinate or an element of $\h^*$ vanishes on its affine cone. This proves the assertion.
\end{proof}

\Needspace{5\baselineskip}
\begin{lemma}\label{lem:basis}
    In simply laced type, for every $d\geq0$, the products in \eqref{eq:chains} form a basis of $\C[\oc\cap\n]_d$.
\end{lemma}

\begin{proof}
    The assertion is immediate for $d=0$. Let $d\geq1$. Expand a homogeneous element of $\C[\oc]$ in the basis $\mathcal B_d$ constructed in the proof of Lemma~\ref{lem:smt}. Every basis monomial containing a negative-root coordinate or an element of the chosen basis of $\h^*$ vanishes in \eqref{eq:quotient}. Every remaining product has the form \eqref{eq:chains}. This proves spanning.

    Suppose that a linear combination of chain products vanishes in the quotient. Fix one term $x_{\gamma_1}\cdots x_{\gamma_d}$ and restrict to the affine cone over $X_{\gamma_1}$. All linear equations in \eqref{eq:quotient} vanish on this cone, so the restriction factors through $\C[\oc\cap\n]$. By Lemmas~\ref{lem:rootorder} and~\ref{lem:smt}, the surviving chain products restrict to distinct basis elements, and the chosen term is among them. Its coefficient is therefore zero. Applying this to each term proves independence.
\end{proof}

Write $I_2=I\cap\C[\n]_2$. For $i<j$ with $\beta_i,\beta_j$ incomparable, expand the class of $x_{\beta_i}x_{\beta_j}$ in the chain basis of Lemma~\ref{lem:basis}. This gives a unique relation
\begin{equation}\label{eq:straighten}
    g_{ij}=x_{\beta_i} x_{\beta_j}-
    \sum_{\substack{k\leq l,\ \beta_k,\beta_l\text{ comparable}\\
    \beta_k+\beta_l=\beta_i+\beta_j}}c_{ij}^{kl}x_{\beta_k}x_{\beta_l}\in I_2.
\end{equation}
The weight condition follows from $T$-equivariance.

\begin{lemma}\label{lem:leading}
    Every nonzero coefficient in \eqref{eq:straighten} satisfies $k<\min\{i,j\}$. Consequently,
    $$
        \LT_\prec(g_{ij})=x_{\beta_i} x_{\beta_j}.
    $$
\end{lemma}

\begin{proof}
    Fix a term with $c_{ij}^{kl}\ne0$. Since $k\leq l$ and $\beta_k,\beta_l$ are comparable, one has $\beta_k\leq\beta_l$. By Lemma~\ref{lem:smt}, their product restricts to a basis element on the affine cone over $X_{\beta_k}$. Restrict \eqref{eq:straighten} to this cone. If either $x_{\beta_i}$ or $x_{\beta_j}$ vanished there, the left product would restrict to zero. By Lemmas~\ref{lem:rootorder} and~\ref{lem:smt}, the surviving products on the right restrict to distinct basis elements, including $x_{\beta_k}x_{\beta_l}$. Their linear independence would force $c_{ij}^{kl}=0$. Thus both coordinates restrict nontrivially. By the coordinate criterion in Lemma~\ref{lem:rootorder},
    $$
        \beta_k\leq\beta_i,\qquad\beta_k\leq\beta_j.
    $$
    Both inequalities are strict, since $\beta_i,\beta_j$ are incomparable. Hence
    $$
        \height(\beta_k)<\min\{\height(\beta_i),\height(\beta_j)\},
        \qquad k<\min\{i,j\}.
    $$

    The least variable at which the exponent vectors differ is $x_{\beta_k}$. Its exponent is zero in $x_{\beta_i}x_{\beta_j}$ and positive in $x_{\beta_k}x_{\beta_l}$. In graded reverse lexicographic order, this means
    \[
        x_{\beta_i}x_{\beta_j}\succ x_{\beta_k}x_{\beta_l}.\qedhere
    \]
\end{proof}

\begin{theorem}\label{thm:groebner}
    In simply laced type, the polynomials $g_{ij}$ form the reduced Gr\"obner basis of $I$, with
    $$
        \inid_\prec I=(x_\alpha x_\beta\mid\alpha,\beta\in\Phi^+\text{ are incomparable}).
    $$
    In particular, $I$ is radical.
\end{theorem}

\begin{proof}
    Divide $f\in I$ by the $g_{ij}$. By Lemma~\ref{lem:leading}, no remainder monomial is divisible by a product of incomparable root coordinates. The remainder is therefore a linear combination of chain monomials. It still belongs to $I$ because every $g_{ij}$ does. Lemma~\ref{lem:basis} forces it to be zero. Thus every element of $I$ reduces to zero, proving the Gr\"obner basis assertion and the formula for the initial ideal.

    The $g_{ij}$ are monic, and their remaining monomials have chain support. None is divisible by a leading monomial, so the basis is reduced.

    If $f^m\in I$, replace $f$ by its remainder. For a nonzero remainder, squarefreeness of the initial ideal would imply
    $$
        \LT_\prec(f)^m\in\inid_\prec I
        \quad\Longrightarrow\quad
        \LT_\prec(f)\in\inid_\prec I.
    $$
    This contradicts the definition of a remainder. Thus $f\in I$.
\end{proof}

The basis can be computed from the restricted quadrics in \eqref{eq:restricted}. Order the degree-two monomials decreasingly under $\prec$ and put the coefficient matrix in reduced row echelon form. Its nonzero rows are the $g_{ij}$. In particular,
$$
    \dim I_2=\#\{\{\alpha,\beta\}\subset\Phi^+\mid
    \alpha,\beta\text{ are incomparable}\}.
$$

\begin{proposition}\label{prop:degeneration}
    Assume that $\Phi$ is simply laced. Let $\Delta(\Phi^+)$ be the simplicial complex whose faces are the chains in the positive-root poset, including the empty chain. There is a flat Gr\"obner degeneration of $\C[\oc\cap\n]$ to $\C[\Delta(\Phi^+)]$. Moreover,
    \begin{align*}
        \Hilb(\C[\oc\cap\n],t)&=\sum_{C\in\Delta(\Phi^+)}
        \left(\frac{t}{1-t}\right)^{|C|},
        \\
        \dim\C[\oc\cap\n]&=\height(\theta).
    \end{align*}
\end{proposition}

\begin{proof}
    The minimal nonfaces of $\Delta(\Phi^+)$ are the incomparable pairs, so its Stanley--Reisner ideal is $\inid_\prec I$. The standard Gr\"obner deformation \cite[\S22]{Pee11} has the stated special fiber. Each variable in the support $C$ of a basis monomial has positive exponent. The contribution of this support to the Hilbert series is therefore
    $$
        \prod_{\alpha\in C}(t+t^2+\cdots)
        =\left(\frac{t}{1-t}\right)^{|C|}.
    $$

    Every positive root is below $\theta$ and above a simple root. Lemma~\ref{lem:rootorder} inserts all intermediate heights, so every maximal chain has $\height(\theta)$ elements. This is the dimension of the Stanley--Reisner ring.
\end{proof}

\Needspace{7\baselineskip}
\begin{lemma}\label{lem:minimal-cm}
    In simply laced type, $\C[\Delta(\Phi^+)]$ is Cohen--Macaulay. Consequently, $\oc\cap\n$ is Cohen--Macaulay.
\end{lemma}

\begin{proof}
    For a root interval $[\gamma,\eta]$, let $\Delta([\gamma,\eta])$ be its complex of chains. Every such interval is totally semimodular. This means that in every subinterval, two distinct elements covering the same element have a common cover. Suppose that $\gamma+\alpha_i$ and $\gamma+\alpha_j$ cover $\gamma$ and lie below $\eta$, with $i\ne j$. Lemma~\ref{lem:rootorder} shows that all covers have this form. Normalize the root lengths to have square $2$. Then
    \begin{align*}
        (\gamma,\alpha_i)&=(\gamma,\alpha_j)=-1,\\
        \|\gamma+\alpha_i+\alpha_j\|^2&=2+2(\alpha_i,\alpha_j).
    \end{align*}
    If $(\alpha_i,\alpha_j)=-1$, the last expression would be zero, although $\gamma+\alpha_i+\alpha_j$ is a nonzero positive combination of simple roots. Hence $(\alpha_i,\alpha_j)=0$, so $(\gamma+\alpha_i,\alpha_j)=-1$. The root-string property now implies that $\gamma+\alpha_i+\alpha_j$ is a root. It covers both given roots and is at most $\eta$, by comparison of simple-root coefficients.

    By \cite[Corollary~5.2]{BW83}, these intervals are shellable. Their chain complexes are pure, with one vertex at each height in a maximal chain. Pure shellable complexes are Cohen--Macaulay \cite[\S5]{Sta77}. Thus for every $\gamma\in\Phi^+$, the Stanley--Reisner ring $\C[\Delta([\gamma,\theta])]$ is Cohen--Macaulay, with
    $$
        \dim\C[\Delta([\gamma,\theta])]
        =\height(\theta)-\height(\gamma)+1.
    $$

    Order the simple roots so that every initial segment spans a connected subtree of the Dynkin diagram. If $\beta$ is not the first root, let $\delta$ be its unique neighbor among the earlier vertices. Their sum $\beta+\delta$ is a root, and
    \begin{equation}\label{eq:tree-intersection}
        \Delta([\beta,\theta])\cap
        \bigcup_{\alpha\text{ earlier than }\beta}\Delta([\alpha,\theta])
        =\Delta([\beta+\delta,\theta]).
    \end{equation}
    Indeed, the least root $\eta$ of a nonempty chain on the left contains $\beta$ and an earlier simple root in its support. The support of a positive root is connected. Since the Dynkin diagram is a tree, it contains the path joining these two vertices, hence also $\delta$. Thus $\eta\geq\beta+\delta$. Conversely, a chain on the right lies above both $\beta$ and the earlier root $\delta$.

    For two simplicial complexes $\Lambda$ and $\Gamma$, there is an exact sequence
    $$
        0\longrightarrow\C[\Lambda\cup\Gamma]
        \longrightarrow\C[\Lambda]\oplus\C[\Gamma]
        \longrightarrow\C[\Lambda\cap\Gamma]
        \longrightarrow0.
    $$
    All rings are modules over the polynomial ring on the combined vertex set. The last map is the difference of the restriction maps. By the depth lemma, the union has a Cohen--Macaulay ring of dimension $d$ if the two summands are Cohen--Macaulay of dimension $d$ and the intersection ring is Cohen--Macaulay of dimension $d-1$.

    Apply this exact sequence to the intervals above the simple roots in the chosen order. Each interval has a Cohen--Macaulay ring of dimension $\height(\theta)$, and by \eqref{eq:tree-intersection} its overlap with the preceding union has a Cohen--Macaulay ring of dimension $\height(\theta)-1$. Every positive-root chain lies above a simple root, so induction proves that $\C[\Delta(\Phi^+)]$ is Cohen--Macaulay.

    By Proposition~\ref{prop:degeneration}, the initial quotient is $\C[\Delta(\Phi^+)]$. Projective dimension cannot increase on passage from an initial quotient to the original quotient \cite[Theorem~22.9(2)]{Pee11}. Both have dimension $\height(\theta)$. The Auslander--Buchsbaum formula over the polynomial ring $\C[\n]$ therefore implies
    \begin{align*}
        \operatorname{depth}\C[\oc\cap\n]
        &\geq\operatorname{depth}\C[\Delta(\Phi^+)]\\
        &=\height(\theta)
        =\dim\C[\oc\cap\n].\qedhere
    \end{align*}
\end{proof}

\Needspace{11\baselineskip}
\begin{proof}[Proof of Theorem~\ref{thm:main}]
    Lemma~\ref{lem:generic} proves generic reducedness of $\oc\cap\n$.

    If $\Phi$ is simply laced, Theorem~\ref{thm:groebner} and Lemma~\ref{lem:minimal-cm} prove that $\oc\cap\n$ is reduced and Cohen--Macaulay.

    If $\Phi$ is not simply laced, Lemma~\ref{lem:short} proves that $\oc\cap\n$ is nonreduced. An associated prime of its nonzero nilradical is also an associated prime of its coordinate ring. Generic reducedness excludes minimal primes from the support of the nilradical, so such a prime is embedded. Thus $\oc\cap\n$ does not satisfy $(S_1)$ \cite[Tag~031Q]{Stacks}.
\end{proof}

The Killing form identifies the Springer resolution \cite[\S\S2.1, 2.10]{BM83} with
$$
    \mu:T^*(G/B)=G\times^B\n\longrightarrow\mathcal N,
    \qquad [g,x]\longmapsto\Ad(g)x,
$$
where $\mathcal N$ is the nilpotent cone. Its scheme-theoretic fibers over $\omin$ are called minimal Springer fibers.

\Needspace{10\baselineskip}
\begin{proposition}\label{prop:springer}
    Every minimal scheme-theoretic Springer fiber is generically reduced. If $\Phi$ is simply laced, it is reduced and Cohen--Macaulay. If $\Phi$ is not simply laced, it is nonreduced and, being generically reduced, fails $(S_1)$.
\end{proposition}

\begin{proof}
    Fix $e\in\omin$ and set $U=(\oc\cap\n)\setminus\{0\}$. Since $\oc\setminus\{0\}=\omin$, the scheme $U$ is the scheme-theoretic intersection $\omin\cap\n$. Base change of the orbit map
    $$
        G\longrightarrow\omin,\qquad g\longmapsto\Ad(g^{-1})e
    $$
    gives a diagram
    $$
        U\ \longleftarrow\ Z:=G\times_{\omin}U\ \longrightarrow\ \mu^{-1}(e).
    $$
    The middle scheme parametrizes pairs $(g,x)$ with $\Ad(g)x=e$. The left map $(g,x)\mapsto x$ is a torsor for the centralizer $G_e$. This algebraic group is smooth over $\C$, so the map is smooth and surjective. Its action is
    $$
        h\cdot(g,x)=(hg,x).
    $$
    The right map $(g,x)\mapsto[g,x]$ is the base change of $G\times\n\to G\times^B\n$. It is a smooth surjective $B$-torsor, with action
    $$
        (g,x)\cdot b=(gb,\Ad(b^{-1})x).
    $$

    For either map, the pullback of the nilradical is nilpotent. Its quotient is smooth over the reduced base, hence reduced, so this pullback is the nilradical of the middle scheme. Reducedness therefore ascends and descends along both maps, by faithful flatness. Generic reducedness is also preserved and reflected: smooth maps are open, send component generic points to component generic points, and induce faithfully flat local maps with geometrically regular fibers at those points. Cohen--Macaulayness also ascends along smooth maps and descends along smooth surjective maps \cite[Tag~045J]{Stacks}. Theorem~\ref{thm:main} now proves generic reducedness in every type and reducedness and Cohen--Macaulayness in simply laced type.

    In non-simply laced type, Lemma~\ref{lem:short} gives a nonzero nilpotent in the coordinate ring of a principal open subset of $U$. Its pullback to the middle scheme remains nonzero because the left map is faithfully flat. The identity of the pulled-back nilradicals then implies that $\mu^{-1}(e)$ is nonreduced. Together with generic reducedness, this excludes $(S_1)$.
\end{proof}

\Needspace{9\baselineskip}
\begin{example}[Type $D_4$]
    Use $\g=\mathfrak{so}_8$ preserving the symmetric bilinear form
    $$
        \langle u,v\rangle=u_1v_8+u_2v_7+\cdots+u_8v_1.
    $$
    Its Gram matrix and the defining equation of $\g$ are
    $$
        J=(\delta_{i+j,9})_{1\leq i,j\leq8},\qquad X^{\mathsf T}J+JX=0.
    $$
    Take $\h$ to be the diagonal Cartan subalgebra and $\n$ the strictly upper triangular subalgebra of $\g$. Let $\varepsilon_i$ be the $i$-th diagonal coordinate for $1\leq i\leq4$. Take the simple roots
    \begin{align*}
        \alpha_1&=\varepsilon_1-\varepsilon_2,
        &\alpha_2&=\varepsilon_2-\varepsilon_3,\\
        \alpha_3&=\varepsilon_3-\varepsilon_4,
        &\alpha_4&=\varepsilon_3+\varepsilon_4.
    \end{align*}
    After $\beta_i=\alpha_i$ for $1\leq i\leq4$, the remaining positive roots are
    \begin{align*}
        \beta_5&=\alpha_1+\alpha_2,
        &\beta_6&=\alpha_2+\alpha_3,\\
        \beta_7&=\alpha_2+\alpha_4,
        &\beta_8&=\alpha_1+\alpha_2+\alpha_3,\\
        \beta_9&=\alpha_1+\alpha_2+\alpha_4,
        &\beta_{10}&=\alpha_2+\alpha_3+\alpha_4,\\
        \beta_{11}&=\alpha_1+\alpha_2+\alpha_3+\alpha_4,
        &\beta_{12}&=\alpha_1+2\alpha_2+\alpha_3+\alpha_4.
    \end{align*}
    The corresponding root coordinates give
    $$
        X=\begin{pmatrix}
            0&x_1&x_5&x_8&x_9&x_{11}&x_{12}&0\\
            0&0&x_2&x_6&x_7&x_{10}&0&-x_{12}\\
            0&0&0&x_3&x_4&0&-x_{10}&-x_{11}\\
            0&0&0&0&0&-x_4&-x_7&-x_9\\
            0&0&0&0&0&-x_3&-x_6&-x_8\\
            0&0&0&0&0&0&-x_2&-x_5\\
            0&0&0&0&0&0&0&-x_1\\
            0&0&0&0&0&0&0&0
        \end{pmatrix}.
    $$
    For graded reverse lexicographic order with $x_1\prec\cdots\prec x_{12}$, the reduced Gr\"obner basis consists of the following nine nonmonomial quadrics,
    \begin{align*}
        &x_9x_{10}-x_7x_{11}+x_4x_{12},
        &&x_8x_{10}-x_6x_{11}+x_3x_{12},\\
        &x_8x_9+x_5x_{11}+x_1x_{12},
        &&x_7x_8+x_2x_{11},\\
        &x_6x_9+x_2x_{11},
        &&x_6x_7+x_2x_{10},\\
        &x_5x_{10}-x_2x_{11},
        &&x_5x_7-x_2x_9,\\
        &x_5x_6-x_2x_8,
    \end{align*}
    and the following fifteen monomials,
    \begin{align*}
        &x_1x_2,\ x_1x_3,\ x_1x_4,\ x_1x_6,\ x_1x_7,\\
        &x_1x_{10},\ x_2x_3,\ x_2x_4,\ x_3x_4,\ x_3x_5,\\
        &x_3x_7,\ x_3x_9,\ x_4x_5,\ x_4x_6,\ x_4x_8.
    \end{align*}
    In each polynomial, the first term is the leading term.

    The matrix $XJ$ is skew-symmetric. The first displayed quadric is its Pfaffian on rows and columns $1,2,3,4$. Row reduction of all four-by-four Pfaffians and the entries of $X^2$ produces the displayed relations. Setting $x_{12}=1$ and all other $x_i=0$ gives a highest-root vector of rank two and square zero, so these quadrics vanish on $\oc$. Their leading terms are distinct and correspond to the $24$ incomparable positive-root pairs. By Lemma~\ref{lem:basis}, $\dim I_2=24$, so the relations span $I_2$. Their remaining terms have chain support, and they are therefore the reduced basis of Theorem~\ref{thm:groebner}.

    The numbers of chains with $0,1,2,3,4,5$ distinct roots are $1,12,42,64,45,12$, respectively. Proposition~\ref{prop:degeneration} then implies
    $$
        \Hilb(\C[\oc\cap\n],t)=\frac{1+7t+4t^2}{(1-t)^5}.
    $$
\end{example}

\Needspace{12\baselineskip}
\begin{example}[Type $C_2$]
    Use the symplectic form with matrix
    $$
        J=\begin{pmatrix}0&I_2\\-I_2&0\end{pmatrix}
    $$
    and the Cartan subalgebra of matrices $\operatorname{diag}(h_1,h_2,-h_1,-h_2)$. Let $\varepsilon_i(h)=h_i$. Take the simple roots and the remaining positive roots in the order
    \begin{align*}
        \beta_1=\alpha_1&=\varepsilon_1-\varepsilon_2,
        &\beta_2=\alpha_2&=2\varepsilon_2,\\
        \beta_3&=\alpha_1+\alpha_2,
        &\beta_4&=2\alpha_1+\alpha_2.
    \end{align*}
    The corresponding Borel nilradical consists of matrices
    $$
        X=\begin{pmatrix}
            0&x_1&x_4&x_3\\
            0&0&x_3&x_2\\
            0&0&0&0\\
            0&0&-x_1&0
        \end{pmatrix}.
    $$
    Multiplication by $J$ identifies $\mathfrak{sp}_4$ with symmetric matrices. Since the symplectic group is transitive on nonzero vectors, the nonzero matrices $vv^{\mathsf T}$ form the highest-root orbit under the resulting congruence action. Thus $\oc$ is the rank-at-most-one symmetric cone. Its ideal is generated by the two-by-two minors: under $y_{ij}\mapsto v_iv_j$, these minors identify all pairings of each fixed multiset of indices, and hence generate the kernel. Restricting the minors, we obtain
    $$
        I=(x_1^2,x_1x_2,x_1x_3,x_3^2-x_2x_4).
    $$
    These four quadrics form the reduced Gr\"obner basis for graded reverse lexicographic order with $x_1\prec x_2\prec x_3\prec x_4$. The primary decomposition is
    $$
        I=(x_1,x_3^2-x_2x_4)\cap(x_1^2,x_2,x_3).
    $$
    Indeed, write an element of the first prime ideal as $x_1f+(x_3^2-x_2x_4)g$. Membership in $(x_1^2,x_2,x_3)$ forces $f\in(x_1,x_2,x_3)$, which proves the equality. The second ideal is $(x_1,x_2,x_3)$-primary. The nilradical is generated by $x_1$, with
    $$
        \operatorname{Ann}(x_1)=(x_1,x_2,x_3).
    $$
    Its support is the line $x_1=x_2=x_3=0$.

    The dense open set $x_2\ne0$ and the open set $x_4\ne0$ have coordinate rings
    \begin{align*}
        \C[\oc\cap\n][x_2^{-1}]&=\C[x_2,x_2^{-1},x_3],\\
        \C[\oc\cap\n][x_4^{-1}]&=\C[x_4,x_4^{-1},x_1,x_3]/(x_1^2,x_1x_3).
    \end{align*}
    In the second ring, we have eliminated $x_2$. This ring is a Laurent polynomial extension in $x_4$ of $\C[x_1,x_3]/(x_1^2,x_1x_3)$. At $x_4=1$ and $x_1=x_3=0$, the element $x_4-1$ is consequently a nonzerodivisor. In the quotient, the nonzero class of $x_1$ is annihilated by the maximal ideal $(x_1,x_3)$. Thus the local ring has dimension two and depth one.
\end{example}

\begin{thebibliography}{Stacks}

    \bibitem[BM83]{BM83}
    W.~Borho and R.~MacPherson,
    \href{https://www.numdam.org/item/AST_1983__101-102__23_0/}{\emph{Partial resolutions of nilpotent varieties}},
    Ast\'erisque \textbf{101--102} (1983), 23--74.

    \bibitem[Bou02]{Bou02}
    N.~Bourbaki,
    \emph{Lie Groups and Lie Algebras, Chapters 4--6},
    Elements of Mathematics, Springer-Verlag, Berlin, 2002.

    \bibitem[BW83]{BW83}
    A.~Bj\"orner and M.~Wachs,
    \href{https://doi.org/10.1090/S0002-9947-1983-0690055-6}{\emph{On lexicographically shellable posets}},
    Trans. Amer. Math. Soc. \textbf{277} (1983), no.~1, 323--341.

    \bibitem[Gar82]{Gar82}
    D.~Garfinkle,
    \href{https://hdl.handle.net/1721.1/15620}{\emph{A new construction of the Joseph ideal}},
    Ph.D. thesis, Massachusetts Institute of Technology, 1982.

    \bibitem[JL25]{JL25}
    B.~Jia and Y.~Li,
    \emph{Minimal nilpotent orbits and toric varieties},
    preprint, 2025,
    \href{https://arxiv.org/abs/2511.02179}{arXiv:2511.02179}.

    \bibitem[Jos98]{Jos98}
    A.~Joseph,
    \href{https://doi.org/10.1016/S0012-9593(98)80017-7}{\emph{Orbital varieties of the minimal orbit}},
    Ann. Sci. \'{E}cole Norm. Sup. (4) \textbf{31} (1998), no.~1, 17--45.

    \bibitem[Lit94]{Lit94}
    P.~Littelmann,
    \href{https://doi.org/10.1007/BF01231564}{\emph{A Littlewood--Richardson rule for symmetrizable Kac--Moody algebras}},
    Invent. Math. \textbf{116} (1994), 329--346.

    \bibitem[Lit96]{Lit96}
    P.~Littelmann,
    \href{https://doi.org/10.1006/aima.1996.0085}{\emph{A plactic algebra for semisimple Lie algebras}},
    Adv. Math. \textbf{124} (1996), no.~2, 312--331.

    \bibitem[Lit98]{Lit98}
    P.~Littelmann,
    \href{https://doi.org/10.1090/S0894-0347-98-00268-9}{\emph{Contracting modules and standard monomial theory for symmetrizable Kac--Moody algebras}},
    J. Amer. Math. Soc. \textbf{11} (1998), no.~3, 551--567.

    \bibitem[Pee11]{Pee11}
    I.~Peeva,
    \href{https://doi.org/10.1007/978-0-85729-177-6}{\emph{Graded Syzygies}},
    Algebra and Applications \textbf{14}, Springer, London, 2011.

    \bibitem[Sta77]{Sta77}
    R.~P.~Stanley,
    \href{https://doi.org/10.1007/978-94-010-1220-1_3}{\emph{Cohen--Macaulay complexes}},
    in \emph{Higher Combinatorics} (M.~Aigner, ed.),
    NATO Adv. Study Inst. Ser. C \textbf{31}, D.~Reidel, Dordrecht, 1977, 51--62.

    \bibitem[Stacks]{Stacks}
    The Stacks Project Authors,
    \emph{The Stacks Project},
    \url{https://stacks.math.columbia.edu}, accessed 25 October 2026.

\end{thebibliography}

\end{document}
